\scanpage{84}
$\mathfrak S_3$ opère (\unsure{fid\textsuperscript{t}|fidt}) sur $\deleted{\unsure{\Pi}}\Omega_0$ \unsure{plutôt|puis} $\Pi$ par

\[
 \sigma(h_i)=\begin{cases}
 h_{\sigma(i)}&\text{si }\sigma\in\mathfrak S_3^+\\
 h_{\sigma(i)}^{-1}&\text{si }\sigma\in\mathfrak S_3^-
 \end{cases}
\]

On a donc des injections

\diagram{\begin{tikzcd}[ampersand replacement=\&,row sep=small]
 \mathfrak S_3\arrow[r,hook]\arrow[dr,hook]\&\Omega_0\arrow[d]\\
 \&\Omega
\end{tikzcd}}

d'où

\diagram{\begin{tikzcd}[ampersand replacement=\&,row sep=small]
 \mathfrak S_3\arrow[r,hook]\arrow[dr,hook]\&\widetilde\Omega_0\arrow[d]\\
 \&\widetilde\Omega
\end{tikzcd}}

[l'injectivité $\mathfrak S_3\unsure{-|\longrightarrow}\widetilde\Gamma$ est vérifiée immédiatement]

\[
 \text{Soit}\qquad\gamma_0=\mathbb Z^*,\qquad
 \gamma=\widehat{\mathbb Z}^{*}
\]

On va définir trois op. de $\gamma_0\times\gamma_0$ \deleted{sur} \unsure{ds} $\Omega_0$, et de $\gamma\times\gamma$ sur $\Omega$, \deleted{notées $\rho$} Si $\varepsilon,\varepsilon'\in\gamma_0\times\gamma_0$ ou $\gamma\times\gamma$, on notera par $\rho_{\bullet\varepsilon,\varepsilon'},\ \rho_{\varepsilon\bullet\varepsilon'},\ \rho_{\varepsilon,\varepsilon'\bullet}$ les effets de $(\varepsilon,\varepsilon')$ \deleted{\unsure{p}} par ces trois opérations. On pose

\[
\left\{\begin{aligned}
 \rho_{\varepsilon,\varepsilon'\bullet}(h_0)&=h_0^\varepsilon,&
 \rho_{\varepsilon,\varepsilon'\bullet}(h_1)&=h_1^{\varepsilon'}&
 \inserted{\text{donc}}\quad\rho_{\varepsilon,\varepsilon'}(\deleted{\unsure{t}}h_2)&=h_1^{-\varepsilon'}h_0^{-\varepsilon}\\
 \rho_{\deleted{\unsure{\varepsilon}}\bullet\varepsilon\varepsilon'}(h_1)&=h_1^\varepsilon,&
 \rho_{\bullet\varepsilon\varepsilon'}(h_2)&=h_2^{\varepsilon'}&
 \rho_{\bullet\varepsilon\varepsilon'}(h_0)&=h_2^{-\varepsilon'}h_2^{-\varepsilon}\\
 \rho_{\varepsilon'\bullet\varepsilon}(h_2)&=h_2^\varepsilon,&
 \rho_{\varepsilon'\bullet\varepsilon}(h_0)&=h_0^{\varepsilon'}&
 \rho_{\varepsilon'\bullet\varepsilon}(h_1)&=h_0^{-\varepsilon'}h_2^{-\varepsilon}
\end{aligned}\right.
\]

\deleted{On définit ainsi, si $I$ est une partie à deux él. de $\{0,1,2\}$} De façon plus intrinsèque, si $I$ est une partie à $2$ él. \deleted{\unsure{de}} de $\{0,1,2\}=\Delta_2$, on définit des hom

\diagram{\begin{tikzcd}[ampersand replacement=\&,row sep=tiny,column sep=small]
 \rho_0^I:\&\gamma_0^I\arrow[r]\arrow[d,hook]\&\unsure{\Gamma_0|\Pi_0}\arrow[d,hook]\\
 \rho^I:\&\gamma^I\arrow[r]\&\Gamma
\end{tikzcd}}

en \deleted{\unsure{posant, si}}\inserted{faisant opérer} $\varepsilon=(\varepsilon_i)_{i\in I}$ sur $\Omega_0$ (resp $\Omega$) par \deleted{l'action de $\varepsilon$}

\[
 \rho^I_{\unsure{\bullet|0}}(\varepsilon)(\underline h_i)=\underline h_i^{\varepsilon_i}\qquad(i\in I)
\]

\deleted{et prolonge} ce qui donne un autom de $\Omega_0\simeq\mathbb Z^{(*I)}$ -- puissance libre

Les relations entre les op. de $\mathfrak S_3$ et les $\rho^I$
% archive sheet number: 83
